\chapter{Discussion}
\label{ch:discussion}

\draftnote{this chapter is a placeholder skeleton built around the
draft-note prompts left in
\cref{ch:resolution,ch:tl-horizontal,ch:tl-vertical,ch:phase-plane,ch:fluidflow}.
Resolve those prompts first (they ask for specific numbers read off specific
figures), then rewrite this chapter as connected prose rather than a list.}

\section{Amplitude versus Phase: Quantifying the Improvement}

\Cref{ch:resolution,ch:tl-horizontal,ch:tl-vertical} established an
amplitude-based detectability floor, and \cref{ch:phase-plane} showed the
phase-plane estimator remaining accurate at scales where that amplitude
floor has already been reached. \draftnote{state the improvement factor
explicitly: floor scale found in \cref{sec:res-comparison,sec:tl-detectability,sec:vtl-detectability}
versus the smallest scale at which \cref{sec:tlp-horizontal-validation,sec:tlp-vertical}
still recover $(\Dz,\Dx)$ within tolerance.}

\section{Lateral/Vertical Asymmetry}

\Cref{sec:th-duality} predicted, from first principles, that a migrated
image separates lateral spatial frequency like a prism but not vertical
frequency. \draftnote{state whether this asymmetry was visible only in the
instantaneous-phase analysis (\cref{fig:tlp-instphase-horiz-summary} vs.\
\cref{fig:tlp-instphase-vert-summary}) or also at the amplitude level
(\cref{fig:tl-summary} vs.\ \cref{fig:vtl-summary}), and what that implies
for a field deployment that cares more about one direction than the other.}

\section{Noise Robustness}

Direct evidence for noise robustness in this thesis is limited to the
lateral, Laplace-noise case of \cref{sec:tl-noise,sec:tlp-noise}. The design
notes behind the WLS estimator (\cref{sec:th-wls}) additionally propose
three further stress tests --- an orthogonality test isolating pure
geometric shift from pure material-change phase, a controlled SNR sweep from
$+30\,\si{dB}$ to $-10\,\si{dB}$ comparing OLS against the WLS estimator, and
a $20\times20$ grid sweep of true $(\Dt,\Dtheta)$ pairs reported as a 2D
inversion-error heatmap --- for which no corresponding output figures were
found under \texttt{TimeLapse\_Figures} at the time this skeleton was
written. \draftnote{either run these three experiments and add the
resulting figures to \cref{ch:phase-plane}, or explicitly scope them as
future work in \cref{ch:conclusion} --- do not leave the gap unstated.}

\section{Limitations}

\draftnote{list concrete limitations: synthetic gprMax data only (no field
validation); back-propagation migration not run on noisy or vertical-noise
data due to computational cost; the domain-size and STFT-scale-labelling
uncertainties flagged in \cref{ch:methodology,sec:tlp-stft}; the fluid-front
model uses a single idealised thin-layer geometry rather than a swept range
of fracture thicknesses or contrasts.}
